\subsection{Phân tích các bên liên quan}

Việc phát triển và triển khai hệ thống mang lại giá trị gia tăng rõ rệt cho các bên liên quan thông qua việc giải quyết trực tiếp các nhu cầu thực tế trong vận hành và kinh doanh. Những nhu cầu này chính là cơ sở để xác lập danh mục các yêu cầu chức năng và phi chức năng chi tiết ở các mục tiếp theo.

\subsubsection{Khách hàng}

Đối với người dùng cuối, hệ thống cung cấp một nền tảng mua sắm hiện đại, tập trung vào tính tiện dụng và cá nhân hóa. Nhu cầu cốt lõi bắt đầu từ các bước thiết lập tài khoản an toàn \textit{[FR1--FR4]} và quản lý thông tin cá nhân thuận tiện \textit{[FR5]}. Trong quá trình mua sắm, khách hàng cần các công cụ tìm kiếm, lọc sản phẩm minh bạch \textit{[FR6, FR9]} và quy trình giỏ hàng, thanh toán đa phương thức linh hoạt \textit{[FR7, FR8]}. 

Điểm nhấn quan trọng nhất là sự hỗ trợ thông minh từ trợ lý AI trong suốt vòng đời mua hàng, từ tư vấn chọn sản phẩm đến hỗ trợ sau mua \textit{[FR11--FR13]}. Để thỏa mãn nhóm người dùng này, hệ thống phải đáp ứng các tiêu chuẩn khắt khe về hiệu năng phản hồi \textit{[NFR1, NFR2]}, tính đơn giản trong thao tác \textit{[NFR3, NFR5]}, cùng khả năng giao tiếp tự nhiên và dẫn dắt người dùng hiệu quả của AI Agent \textit{[NFR4, NFR6]}.

\subsubsection{Đội ngũ vận hành và quản trị viên}

Đối với đội ngũ nội bộ, hệ thống đóng vai trò là công cụ tối ưu hóa hiệu suất làm việc. Nhu cầu chính là khả năng kiểm soát quyền hạn truy cập \textit{[FR16]} và quản lý danh mục sản phẩm, khuyến mãi một cách tập trung \textit{[FR17]}. Quản trị viên cần được hỗ trợ theo dõi lịch sử, đánh giá từ khách hàng \textit{[FR14, FR15]} và cập nhật trạng thái đơn hàng một cách chính xác \textit{[FR19]}. 

Về phía kỹ thuật và vận hành, nhóm này yêu cầu hệ thống phải duy trì độ ổn định cao để phục vụ demo và vận hành thực tế \textit{[NFR7]}. Đặc biệt, các báo cáo quản trị phải được hiển thị nhanh chóng \textit{[NFR10]} và cơ chế đồng bộ dữ liệu phải đảm bảo tính toàn vẹn tuyệt đối \textit{[NFR11]}, tránh các sai sót về logic tồn kho hoặc dòng dữ liệu giữa các dịch vụ.

\subsubsection{Chủ sở hữu doanh nghiệp}

Về mặt kinh doanh, hệ thống được kỳ vọng sẽ nâng cao hiệu quả tiếp cận thị trường và tối ưu hóa doanh thu thông qua các báo cáo thống kê trực quan \textit{[FR18]}. Việc cho phép khách hàng thực hiện đầy đủ các quyền lợi như hủy đơn hàng theo ràng buộc \textit{[FR10]} giúp xây dựng niềm tin và sự chuyên nghiệp cho doanh nghiệp. 

Ngoài ra, hệ thống còn đáp ứng đòi hỏi đảm bảo an toàn cho tài sản dữ liệu và thông tin khách hàng thông qua các cơ chế bảo mật nghiêm ngặt \textit{[NFR8]}. Đồng thời, kiến trúc hệ thống phải có tính mở, đảm bảo khả năng mở rộng linh hoạt \textit{[NFR9]} để sẵn sàng tích hợp thêm các dịch vụ giá trị gia tăng trong tương lai mà không cần tái cấu trúc nền tảng.

\subsubsection{Nhóm phát triển đề tài}

Đồ án là môi trường thực chứng để nhóm ứng dụng các kỹ thuật tiên tiến về vi dịch vụ và trí tuệ nhân tạo. Nhu cầu của nhóm là xây dựng một hệ thống có tính module hóa cao, đảm bảo luồng sự kiện được truyền tải thông suốt giữa các Microservices cốt lõi. Việc đáp ứng đầy đủ các yêu cầu chức năng từ \textit{[FR1]} đến \textit{[FR19]} và các tiêu chuẩn phi chức năng từ \textit{[NFR1]} đến \textit{[NFR11]} là minh chứng cho khả năng tích hợp công nghệ AI vào một hệ thống thương mại điện tử thực tế, tạo ra sản phẩm có giá trị nghiên cứu và ứng dụng cao.